\documentclass[12pt,a4paper]{report}
\usepackage[utf8]{inputenc}
\usepackage[T2A]{fontenc} % T2A лучше подходит для кириллицы
\usepackage[russian]{babel}
\usepackage{amsmath,amssymb}
\usepackage{geometry}
\usepackage{booktabs}
\usepackage{hyperref}
\geometry{margin=2.5cm}
\setlength{\parindent}{0pt}
\setlength{\parskip}{0.8em}

\title{\textbf{Пока учёные ищут окружности: \\ Парадигма IT3 и объединение космологии}\\[0.5cm]
\large Геометрический и термодинамический синтез}
\author{Виктор Логвинович\\
Магистр физико-математических наук\\
Независимый исследователь в области космологических наук\\
Кобрин, Беларусь}
\date{10 апреля 2026 г.}

\begin{document}

\maketitle
\tableofcontents
\newpage

\chapter*{Предисловие: пока учёные ищут окружности}
Более двух десятилетий космологическое сообщество ждёт неопровержимого доказательства конечности Вселенной: парных совпадающих окружностей в реликтовом излучении. Теория элегантна, математика безупречна, но окружности остаются ненайденными. Пока этот поиск продолжается, глубже уже решён другой вопрос: \textit{нужна ли Вселенной бесконечность, чтобы работать?}

Парадигма плоского иррационального тора (IT3) показывает, что нет. Рассматривая Вселенную как конечное многократно связное пространство с определённой геометрической формой, мы можем объяснить аномалию низких мультиполей в реликтовом излучении, снять напряжение Хаббла, заменить тёмную энергию физическим механизмом диссипации и обеспечить бесконечную термодинамическую стабильность. Окружности — прекрасный тест, но не фундамент. Фундамент — это геометрия, поток энергии и крупномасштабная структура. В этой книге мы шаг за шагом разберём этот фундамент простым языком, показывая, как конечная Вселенная не просто выживает, но и развивается.

\chapter{Форма пространства: конечная Вселенная}
\section{Почему конечная?}
Стандартная космология предполагает, что пространство бесконечно и односвязно. Это допущение удобно для расчётов, но оставляет две устойчивые загадки без ответа:
\begin{enumerate}
    \item \textbf{Аномалия низких мультиполей:} Реликтовое излучение показывает меньшие температурные флуктуации на самых больших угловых масштабах, чем ожидается. В бесконечной Вселенной должны существовать волны любой длины. В конечной волны, превышающие размер самого пространства, просто не помещаются.
    \item \textbf{Напряжение Хаббла:} Измерения ранней Вселенной дают $H_0 \approx 67.4$ км/с/Мпк, тогда как локальные измерения дают $H_0 \approx 73.0$ км/с/Мпк. Если история расширения Вселенной немного отличается от стандартной, оба значения могут быть верны.
\end{enumerate}

\section{Аналогия с тором}
Представьте карту из классической видеоигры. Если вы уходите за правый край, вы появляетесь слева. Если вы улетаете за верхний край, вы появляетесь снизу. Это \textbf{тор}. Он конечен, но у него нет краёв. Свет, летящий в одном направлении, в итоге возвращается с противоположной стороны.

Модель IT3 использует трёхмерный тор с длинами сторон $(L_x, L_y, L_z)$. Вместо произвольных соотношений мы фиксируем их иррациональными числами:
$$
\frac{L_y}{L_x} = \sqrt{2}, \qquad \frac{L_z}{L_x} = \sqrt{3}.
$$
Зачем иррациональные? Потому что рациональные соотношения создают повторяющиеся паттерны, которые могут искусственно выстраивать структуры. Иррациональные соотношения гарантируют, что траектории в пространстве никогда не совпадут точно, имитируя наблюдаемую случайность, сохраняя при этом конечность пространства.

\section{Что говорят данные}
Когда мы сопоставляем эту топологию с данными Planck, математика строго ограничивает фундаментальный масштаб:
$$
L_x = 28.57^{+0.73}_{-0.87}~\mathrm{Гпк}.
$$
Это не произвольное число. Оно находится ровно на геометрическом пределе: Вселенная достаточно велика, чтобы вместить наблюдаемый горизонт, но достаточно мала, чтобы подавить самые крупные флуктуации реликтового излучения. Модель улучшает согласие с дефицитом квадруполя на $\Delta\chi^2 = -5.33$ и естественно сдвигает предпочтительную постоянную Хаббла вверх, снижая напряжение с локальными измерениями.

Простыми словами: конечная форма отсекает самые длинные волны, и это отсечение вынуждает скорость расширения адаптироваться. Геометрия объясняет то, для чего в стандартной модели требуются новые частицы или тонкая настройка.

\chapter{Термодинамический двигатель: энергия без тёмной энергии}
\section{Проблема конечной Вселенной}
Если Вселенная конечна и заполнена материей, гравитация должна в конечном итоге стянуть всё обратно. Замкнутая Вселенная коллапсирует. Открытая расширяется вечно, но остывает до абсолютного нуля. Стандартная космология обходит это, добавляя космологическую постоянную $\Lambda$, но у $\Lambda$ нет физического механизма, и она приводит к печально известному расхождению вакуумной энергии на $10^{120}$ порядков.

Конечная топология предлагает другой путь: \textbf{дренаж энергии}. Если пространство компактно, энергия, высвобождаемая при формировании структур (слияния галактик, активность чёрных дыр, гравитационные волны), не рассеивается бесконечно. Она интерферирует сама с собой, образуя стоячие волны, которые фокусируются к симметричным точкам в фундаментальной области.

\section{Механизм стока}
Мы моделируем эту фокусировку как макроскопическое отрицательное давление:
$$
P_{\mathrm{sink}} = -\kappa \rho_m^2, \qquad \kappa = \eta G L_x^2.
$$
Здесь $\rho_m$ — плотность материи, $G$ — гравитационная постоянная, $L_x$ — масштаб топологии, а $\eta$ — геометрический коэффициент эффективности. Квадратичная зависимость ($\rho^2$) возникает потому, что инжекция энергии масштабируется с взаимодействиями (слияния, столкновения), которые зависят от квадрата плотности.

Это модифицирует стандартный закон сохранения энергии:
$$
\dot{\rho}_m + 3H\rho_m = -\Gamma_{\mathrm{sink}}\rho_m^2, \qquad \Gamma_{\mathrm{sink}} = \eta' \frac{G L_x}{c}.
$$
Простыми словами: по мере расширения Вселенной материя разбавляется обычным образом, но небольшая её часть непрерывно «стекает» в топологический «центр». Этот сток действует как отрицательное давление, расталкивая пространство.

\section{Замена тёмной энергии}
Когда мы подставляем это в уравнения Эйнштейна, член стока генерирует эффективный космологический параметр:
$$
\Lambda_{\mathrm{eff}} = \frac{8\pi G \kappa \langle \rho^2 \rangle}{c^2}.
$$
В отличие от истинной постоянной, $\Lambda_{\mathrm{eff}}$ эволюционирует вместе с плотностью. В ранней Вселенной он ничтожен. По мере формирования структур и падения $\rho_m$ сток постепенно доминирует, обуславливая наблюдаемое ускоренное расширение. Численное интегрирование показывает, что Вселенная избегает как Большого сжатия, так и тепловой смерти. Она переходит в стабильное неравновесное состояние. Поле тёмной энергии не нужно. Только геометрия и диссипация.



\chapter{Космическая губка: как структура усиливает геометрию}
\section{Войды не пустые}
Обзоры галактик показывают, что $\sim 75\%$ объёма космоса занято войдами, соединёнными тонкими филаментами. Это не случайный шум; это перколяционная сеть с фрактальной размерностью $D_f \approx 2.6$. В рамках IT3 эта сеть функциональна. Она действует как волновод с низким сопротивлением для релятивистских носителей энергии.

\section{Фактор каналирования}
Мы количественно оцениваем усиление с помощью теории перколяции:
$$
\xi(\phi) = \left(\frac{\phi}{1-\phi}\right)^\alpha, \qquad \alpha \approx 1.2,
$$
где $\phi$ — доля войдов. Сегодня $\phi_0 \approx 0.75$, что даёт $\xi_0 \approx 3.8$. Это означает, что структура губки умножает эффективность стока почти в четыре раза. Ключевой момент: $\xi(\phi) \to 0$ в ранней однородной Вселенной, поэтому механизм остаётся спящим до тех пор, пока не сформируется крупномасштабная структура.

\section{Позднее включение}
Пористость эволюционирует логистически:
$$
\phi(z) = \frac{\phi_0}{1 + \exp[(z - z_{\mathrm{perc}})/\Delta z]}, \quad z_{\mathrm{perc}} \approx 1.0.
$$
Это создаёт функцию активации, которая включается примерно при $z \lesssim 1$. До этого Вселенная расширяется как стандартная $\Lambda$CDM. После этого губка эффективно каналирует энергию, повышая локальную скорость расширения. Это время не подогнано вручную; это прямое следствие момента перколяции войдов.

\chapter{Объединение напряжений: почему правы обе стороны}
\section{Напряжение Хаббла: новый взгляд}
Зонды ранней Вселенной (CMB) измеряют расширение, экстраполируя данные от $z \approx 1100$ при предположении постоянной $\Lambda$. Они видят базовую скорость: $H_0^{\mathrm{early}} \approx 67.4$ км/с/Мпк. Локальные зонды (сверхновые, цефеиды) измеряют расширение при $z < 0.1$, глубоко внутри режима активной губки. Они видят усиленную скорость: $H_0^{\mathrm{local}} \approx 73.0$ км/с/Мпк.

Обе правы. Они измеряют одну и ту же Вселенную в разных топологических фазах.

\section{Сохранение согласованности с CMB}
Естественный вопрос: если локальная скорость расширения меняется, не ломает ли это акустические пики CMB? Нет. Угловое диаметровое расстояние до поверхности последнего рассеяния — это интеграл:
$$
D_A(z_*) = \frac{c}{1+z_*} \int_0^{z_*} \frac{dz}{H(z)}.
$$
Поскольку губка активируется только при $z \lesssim 1$, интеграл определяется физикой ранней Вселенной, где $H(z)$ совпадает с $\Lambda$CDM. Поздняя модификация вносит менее $1\%$ в полное расстояние. Акустический масштаб остаётся нетронутым. Напряжение исчезает не из-за изменения ранней физики, а благодаря пониманию того, что поздняя физика зависит от фазы структуры.

\section{Напряжение $S_8$}
Тот же механизм стока вводит дополнительное трение в рост возмущений плотности. Модифицированное уравнение роста принимает вид:
$$
\ddot{\delta} + (2H + \Gamma_{\mathrm{sink}}^{\mathrm{eff}}\bar{\rho}_m)\dot{\delta} = [4\pi G\bar{\rho}_m + \dots]\delta.
$$
Дополнительное трение подавляет формирование структур примерно на $10\%$ при $z=0$, приводя предсказанный $S_8$ в согласие с обзорами слабого линзирования. Опять же, никаких новых частиц. Только управляемая топологией диссипация, воздействующая на рост структур.



\chapter{Что мы предсказываем: помимо окружностей}
\section{Совпадающие окружности — лишь один из тестов}
Парадигма IT3 предсказывает парные совпадающие окружности в CMB с угловыми радиусами $\alpha \approx 30^\circ\text{--}60^\circ$. Их обнаружение было бы впечатляющим. Их отсутствие не опровергло бы всю теорию, потому что модель делает множество независимых предсказаний:
\begin{enumerate}
    \item \textbf{Зависящая от красного смещения $H_0(z)$:} Выводимое $H_0$ должно монотонно расти при $z < 1$, отслеживая активацию губки.
    \item \textbf{Усиленный эффект Сакса–Вольфа в сверхвойдах:} Топологический поток энергии должен давать $\Delta T/T \sim 10^{-6}$ в войдах с $R > 100$ Мпк, что примерно вдвое превышает предсказание $\Lambda$CDM.
    \item \textbf{Подавление размера войдов:} Модифицированная поддержка давления должна снизить количество войдов с $R > 50$ Мпк примерно на $15\%$.
    \item \textbf{Анизотропия локального $H_0$:} Выравнивание филаментов должно индуцировать дипольную анизотропию $\sim 1\%$ в локальных измерениях расширения.
    \item \textbf{Подавление тензорных мод:} Топологическое трение должно удерживать тензорно-скалярное отношение $r < 0.01$, что согласуется с текущими пределами BICEP/Keck.
\end{enumerate}

\section{Многовекторная структура теории}
Наука продвигается вперёд не ожиданием одного чудесного наблюдения, а построением сетей проверяемых предсказаний. IT3 предоставляет такую сеть. Если найдут окружности — топология подтверждена. Если в войдах появятся аномалии Сакса–Вольфа — губка подтверждена. Если $H_0(z)$ эволюционирует как предсказано — позднее включение подтверждено. Теория стоит или падает на согласованности целого, а не на обнаружении одного признака.

\chapter{Заключение: геометрия, а не мистика}
Современная космология десятилетиями добавляла невидимые компоненты, чтобы сохранить допущение: пространство бесконечно. Парадигма IT3 задаёт более простой вопрос: \textit{что если пространство конечно, и мы просто забыли учесть, как энергия течёт в замкнутой форме?}

Ответ согласован:
\begin{itemize}
    \item \textbf{Форма} отсекает длинные волны, объясняя аномалию низких мультиполей.
    \item \textbf{Поток} отводит энергию в топологические пучности, заменяя тёмную энергию.
    \item \textbf{Структура} каналирует этот поток через войды, активируя ускорение и подавляя избыточное кластеризование.
\end{itemize}

Никаких новых частиц. Никакой тонкой настройки. Никакой мультивселенной. Только геометрия конечного тора, термодинамика диссипации и наблюдаемое губчатое распределение материи. Совпадающие окружности остаются прекрасной целью для будущих экспериментов с CMB. Но пока учёные ищут их, математика уже работает. Напряжения исчезают. Вселенная стабилизируется. Элементы пазла складываются.

Следующее поколение обзоров — CMB-S4, Euclid, DESI Year 5, LSST — решительно проверит эти предсказания. Если данные совпадут, космология перейдёт от парадигмы невидимых субстанций к парадигме видимой геометрии. Если нет — модель будет уточнена или заменена. Так и работает наука. Но на данный момент теоретический аппарат завершён, самосогласован и готов к тому, чтобы звёзды вынесли свой вердикт.

\begin{thebibliography}{99}

\bibitem[Arnold \& Avez(1968)]{Arnold1968}
Arnold V.~I., Avez A., 1968, Ergodic Problems of Classical Mechanics. Benjamin

\bibitem[Cornish et al.(2004)]{Cornish2004}
Cornish N.~J., et al., 2004, Classical and Quantum Gravity, 21, S67

\bibitem[Logvinovich(2026a)]{Logvinovich2026a}
Logvinovich V., 2026a, Zenodo, DOI:10.5281/zenodo.19431323 (Статья I)

\bibitem[Logvinovich(2026b)]{Logvinovich2026b}
Logvinovich V., 2026b, Zenodo, DOI:10.5281/zenodo.19473353 (Статья II)

\bibitem[Logvinovich(2026c)]{Logvinovich2026c}
Logvinovich V., 2026c, Zenodo, DOI:10.5281/zenodo.19479551 (Статья III)

\bibitem[Logvinovich(2026d)]{Logvinovich2026d}
Logvinovich V., 2026d, Zenodo, DOI:10.5281/zenodo.19489307 (Статья IV)

\bibitem[Pan et al.(2012)]{Pan2012}
Pan D.~C., et al., 2012, MNRAS, 421, 926

\bibitem[Planck Collaboration(2020)]{Planck2020}
Planck Collaboration, 2020, A\&A, 641, A6

\bibitem[Riess et al.(2024)]{Riess2024}
Riess A.~G., et al., 2024, ApJL, 934, L7

\end{thebibliography}

\end{document}